% Copyright (c) 2026 Eugene Joseph M. Ragasa.
% SPDX-License-Identifier: Apache-2.0
% Author: Eugene Joseph M. Ragasa
% Formal proof completed: 2026-10-19
\documentclass[11pt]{article}
\usepackage[T1]{fontenc}
\usepackage{amsmath,amssymb,booktabs,tabularx}
\usepackage[margin=1in]{geometry}
\title{Proof Dossier: Exact Finite-Abelian Fourier Inversion}
\author{Eugene Joseph M. Ragasa}
\date{Formal proof completed: 2026-10-19}
\begin{document}\maketitle
\section{Identity and status}
\begin{tabularx}{\textwidth}{@{}lX@{}}\toprule
Statement ID & \texttt{BG-FOURIER-THM-010} \\
Status & proved; no admissions \\
Lean declaration & \texttt{ProjectKoios.BandGauge.Fourier.}\newline\texttt{finiteFourierReconstruction\_coefficients} \\
Owning module & \texttt{ProjectKoios/BandGauge/Fourier/FiniteAbelianTransform.lean} \\
Axioms & \texttt{propext}, \texttt{Classical.choice}, \texttt{Quot.sound} \\
\bottomrule\end{tabularx}
\section{Statement}
Let $f:G\to V$, where $G$ is a finite abelian group and $V$ is any complex module. Define
\[
 \widehat f(\chi)=\frac1{|G|}\sum_{x\in G}\overline{\chi(x)}f(x).
\]
Then $\sum_{\chi\in\widehat G}\chi(k)\widehat f(\chi)=f(k)$ for every $k\in G$.
\section{Argument}
Distribute the finite sums and interchange their order. The scalar multiplying each $f(x)$ is the character kernel from BG-FOURIER-THM-009, so every term vanishes except $x=k$.
\section{Boundary and provenance}
The convention is exact and explicit but is not yet bound to a publication mesh. This theorem is not Parseval, a continuum inversion theorem, or a numerical validation result. Lean checked the exact statement; the dossier and proof were AI-assisted.
\end{document}
